\scanpage{112}

où $\widetilde X,\widetilde Y$ sont les rev. universels de $X,Y$ déduits de \inserted{$T\longrightarrow X$ et de $T\longrightarrow Y$,} de telle façon que l'on obtient
\[
(17)\qquad \mathcal E_{X/Y}^T\simeq\mathcal E_{X/Y}^{\widetilde X}\simeq\operatorname{Ker}\bigl(\mathcal E_X^{\widetilde X}\unsure{\mathrel{\text{---}}|\longrightarrow}\mathcal E_Y^{\widetilde Y}\bigr)
\]
(sans hyp. de suite exacte)

et
\[
\begin{aligned}
(18)\qquad \mathcal E_{X/Y}^T\simeq{}&\deleted{\operatorname{Ker}\bigl(\unsure{\mathcal E_{X_{\widetilde Y}}^{\widetilde X}}\unsure{\mathrel{\text{---}}|\longrightarrow}\unsure{\mathcal E_{\widetilde Y}^{\widetilde Y}}\bigr)}\\
&\inserted{\mathcal E_{X_{\widetilde Y}}^Z\unsure{\simeq}\mathcal E_{X_{\widetilde Y}}^{\widetilde X}}
\end{aligned}
\]
dans le cas général.

\diagram{\begin{tikzcd}
X\arrow[d] & X_{\widetilde Y}\arrow[l]\arrow[d] & \widetilde X\arrow[l] & Z\arrow[l]\\
Y'\arrow[d] & \widetilde Y\arrow[l] & & \\
Y & & &
\end{tikzcd}}

Une autre interprétation de $\mathcal E_{X/Y}^T$ s'obtient en prenant $T$ réduit à un $p^t$ géométrique \inserted{$\unsure{\xi|\zeta}_X$}, qui définit donc un $p^t$ géom. $\unsure{\xi|\zeta}_Y$ de $Y$, une fibre géom. $X_{\unsure{\xi|\zeta}_Y}$ de $X$ sur $Y$ en $\unsure{\xi|\zeta}_Y$, $\unsure{\xi|\zeta}_X$ s'interprétant comme un $p^t$ géom. de cette fibre, et on aura
\[
(19)\qquad \mathcal E_{X/Y}^{\unsure{\xi|\zeta}_X}\simeq\mathcal E_{X_{(\unsure{\xi|\zeta}_Y)}}^{\unsure{\xi|\zeta}_X}\simeq\pi_1(X_{(\unsure{\xi|\zeta}_Y)},\unsure{\xi|\zeta}_X)
\]
--- i.e. $\mathcal E_{X/Y}$ s'interprète comme groupe fondamental des ``fibres géom'' de $X/Y$.

\margin{111}
